\section{Event-Based Vision}
\label{sec:background_eventcamera}

\subsection{Event Cameras}
\label{subsec:bg_event_camera}

\subsubsection*{Frame-Based vs.\ Event-Based Sensing}

Conventional frame-based cameras operate synchronously, where every pixel is read out at a fixed rate, producing a dense intensity matrix regardless of whether anything in the scene has changed since the last frame.
This design introduces several well-documented limitations: motion blur at high speeds, a fixed dynamic range bounded by the exposure window, and substantial data redundancy during static scenes --- all of which are particularly problematic for high-speed automotive perception \cite{gallego2020event}.

An event camera fundamentally departs from this paradigm;
rather than reading out all pixels simultaneously, each pixel operates independently, firing an event the moment its log-intensity changes beyond a threshold.
The result is a continuous, asynchronous stream of discrete events rather than a sequence of synchronous frames, registering only what has changed between scenes instead of the full scenes themselves.
In other words, if the scene was completely dormant with no changes occurring, an event-driven camera would capture this as a blank image versus a frame-based camera capturing the empty scene.

\subsubsection*{Advantages and the Compatibility Problem}

This per-pixel, threshold-driven design gives event cameras a compelling profile for autonomous perception: high temporal resolution, high dynamic range, low latency, and low data redundancy \cite{gallego2020event}.
However, a sparse, unordered point cloud of $(x, y, t, p)$ tuples is fundamentally incompatible with conventional deep learning frameworks, which expect dense, structured tensors of fixed spatial dimensions.
Correcting this to adapt to deep learning approaches requires a preprocessing step to convert the asynchronous stream into a format suitable for neural network ingestion, the specific design of which is detailed in Section~\ref{sec:data_preprocessing}.

\subsection{The Natural SNN--Event Camera Pairing}
\label{subsec:bg_snn_event_pairing}

\subsubsection*{An Event-Driven Architecture for an Event-Driven Sensor}

Both systems are fundamentally event-driven: an event camera outputs discrete binary events triggered by scene changes, whilst an \gls{snn} processes discrete binary events in the form of spikes triggered by membrane threshold crossings.
In both cases, silence is the default operating state, and meaningful activity is generated only in response to genuine stimuli.

A \gls{cnn} must collapse the sparse, asynchronous event stream into a dense synchronous tensor before it can process it, discarding the temporal structure encoded in the event timestamps in the process.
An \gls{snn}, by contrast, processes the preprocessed spatiotemporal tensor iteratively across $T$ timesteps via its \gls{lif} dynamics, where each timestep exposes the network to a temporal slice of the event stream, which the \gls{lif} neurons integrate progressively into their membrane potentials.
This directly mirrors the camera's temporal output structure, preserving and exploiting the event-driven dynamics rather than collapsing them into a single static representation \cite{gallego2020event, pfeiffer_snn_opportunities}.

\subsubsection*{Implications for Neuromorphic Deployment}

On neuromorphic hardware such as Intel Loihi \cite{davies2018loihi} or IBM TrueNorth \cite{merolla2014million}, this alignment extends to the physical level.
An event camera connected directly to a neuromorphic chip can route events to spiking neurons without any voxelization or dense tensor conversion, since the raw events can be routed directly to spiking neurons without any intermediate tensor representation, enabling a fully asynchronous end-to-end inference pipeline and eliminating the preprocessing bottleneck entirely.
